\section{Introduction}
\label{sec:intro}

Reinforcement learning (RL) has become a central post-training tool for improving
the reasoning and alignment behavior of large language models
(LLMs)~\citep{ouyang2022training,guo2025deepseekr1,grpo}. 
In this setting, an LLM is optimized as an autoregressive token-level policy
against scalar sequence-level rewards.
In practice, however, the data used for optimization is not generated by exactly
the same policy being updated.
Rollouts are often produced by inference engines whose numerical behavior differs
from the training stack, creating a training-inference mismatch~\citep{qi2025defeating,yao2025offpolicy}.
Policy staleness adds another mismatch, since the policy is typically updated multiple
times on minibatches drawn from a fixed batch of rollout data~\citep{zheng2025stabilizing}.
Together, these effects make practical LLM RL inherently \emph{off-policy}.

Widely used methods~\citep{ppo,grpo,yu2025dapo,cispo,drgrpo} therefore follow a
common trust-region recipe of optimizing a clipped or masked surrogate objective
on rollouts collected by the behavior policy.
Such masks are \emph{asymmetric}.
Being outside the trust region is not sufficient for masking.
A token is masked only when its advantage-weighted update would further increase
the relevant deviation from the behavior policy; updates that reduce this
deviation remain active.
Thus a mask has two components:
the \emph{proximity} criterion asks whether the trust region is exceeded, and
the \emph{direction} criterion asks whether the update points further outward.
In Proximal Policy Optimization (PPO)~\citep{ppo}, both decisions are made from
the sampled token's importance ratio.
Once the ratio leaves a fixed window around one and the advantage pushes it
further away, the token is masked and its gradient vanishes.

The sampled importance ratio, however, is a poor proxy for distributional shift
in LLMs~\citep{dppo}.
It observes only the sampled token and ignores how probability mass moves over
the rest of the vocabulary.
Recent divergence-mask methods such as DPPO therefore replace the ratio-based
proximity criterion with a divergence between the behavior and training
policies~\citep{dppo}.
This gives a more faithful trust region, typically estimated from the top-$K$
log-probabilities exposed by rollout engines.
This upgrade, however, changes only the proximity criterion.
The direction criterion remains inherited from PPO.
For PPO, this direction criterion is consistent with the proximity criterion,
because both are defined by the sampled importance ratio.
For a divergence mask, however, the proximity criterion is defined by a
distributional divergence.
Whether the next update increases that divergence depends not only on the
sampled token but also on how the whole softmax distribution changes.
Thus DPPO makes the proximity criterion distributional, but leaves the direction
criterion single-sample.

To address the aforementioned inconsistency, we define the direction criterion by asking whether the next gradient
step will \emph{increase the divergence}, and mask a token according to the sign
of this predicted change.
We call the resulting rule the \emph{predictive divergence mask}.
For the softmax policies that every autoregressive LLM uses, this first-order
change has a clean closed form whose coefficient is a simple function of the
current token probabilities.
The coefficient splits into two terms.
A \emph{local} term at the sampled token is exactly the quantity the
ratio-based direction criterion already reads.
A \emph{global} term captures the softmax normalization coupling across the
vocabulary, which a single sampled ratio cannot observe.
Keeping only the local term recovers the ratio-based direction criterion, while
keeping both lets the sign track the true change of the divergence.
In practice, rollout engines usually expose only top-$K$ token probabilities
rather than the full vocabulary distribution.
We therefore adapt the predictive criterion to this truncated view with two
lightweight estimators for the unseen tail.
They correspond to an aggregated-tail approximation and a uniform-tail
approximation, respectively.
Both estimators use only simple reductions over the retained probabilities and
add negligible computational overhead.
Empirically, detailed analysis shows that the direction
criterion used by the predictive divergence mask is better aligned with the
realized change of the divergence than the sampled-ratio direction criterion.
Across model scales and precision settings, the predictive divergence mask
improves training stability and effectiveness over the baselines.
